% !TeX spellcheck = es_ES
\documentclass[letterpaper,11pt]{article}
\usepackage[utf8]{inputenc}
\usepackage[T1]{fontenc}
\usepackage[spanish]{babel}
\decimalpoint % separador decimal con punto, como en las tablas de la fuente original
\usepackage{lmodern,microtype}
\usepackage{amsmath,amssymb,amsfonts}
\usepackage{geometry}
\geometry{margin=2.54cm}
\usepackage{setspace}
\setstretch{1.15}
\usepackage{enumitem}
\usepackage{booktabs}
\usepackage{color}
\usepackage{hyperref}
\definecolor{ntr}{rgb}{0.8,0,0}
\hypersetup{colorlinks=true,linkcolor=ntr,citecolor=ntr,urlcolor=black}

% Portada (Luis Cabral): titulo \LARGE sans-serif alineado a la izquierda, sin numero de pagina
\makeatletter
    \renewcommand{\@maketitle}{%
    \newpage\parindent0in\null\vskip2em%
    {\LARGE\textbf{\textsf{\@title}}\par\vskip2em}%
    \@author\par\vskip2em%
    \@date%
    \par
    \thispagestyle{empty}\setcounter{page}{1}%
    }%
\makeatother

\DeclareMathOperator{\Var}{Var}
\DeclareMathOperator{\Cov}{Cov}
\DeclareMathOperator{\Corr}{Corr}
\DeclareMathOperator{\E}{\mathbb{E}}
\DeclareMathOperator{\N}{Normal}
\DeclareMathOperator{\LP}{\mathbb{L}}
\newcommand\independent{\protect\mathpalette{\protect\independenT}{\perp}}
\def\independenT#1#2{\mathrel{\rlap{$#1#2$}\mkern2mu{#1#2}}}


\title{Análisis Econométrico \\ Ayudantía 3}
\author{Magíster en Economía, Universidad Alberto Hurtado}
\date{}

\begin{document}
\maketitle

\noindent Temas: insesgadez y consistencia, convergencia en probabilidad, ley de los grandes
números, teorema de la función continua, convergencia en distribución, teorema central del
límite, teorema de Slutsky y método delta.

\medskip
\noindent Los ejercicios marcados con \textbf{(*)} son los que alcanzamos a resolver en la
ayudantía (80 minutos); los demás quedan como práctica adicional, con solución en la pauta.

\medskip
\noindent En todos los ejercicios suponga que la muestra $\{y_i\}_{i=1}^N$ (o
$\{(y_i,x_i)\}_{i=1}^N$) es i.i.d.\ y que existen los momentos que se necesiten. En cada
respuesta \textbf{enuncie el teorema que está usando}.

\begin{enumerate}[itemsep=1.2em]

%% ---------------------------------------------------------------- 1
\item \textbf{(*)} \emph{Media muestral ponderada.} La media muestral ponderada se define como
    \begin{align*}
    \bar y^{\,*} = \frac{1}{N}\sum_{i=1}^{N} w_i\, y_i,
    \end{align*}
    donde los $w_i$ son constantes no negativas que satisfacen $\frac{1}{N}\sum_{i=1}^N w_i=1$.
    Sea $\mu=\E(y)$ y $\sigma^2=\Var(y)$.
    \begin{enumerate}
        \item Demuestre que $\bar y^{\,*}$ es un estimador insesgado de $\mu$.
        \item Calcule $\Var(\bar y^{\,*})$.
        \item Encuentre condiciones sobre los $w_i$ para que $\lim_{N\to\infty}\Var(\bar y^{\,*})=0$,
              y dé un ejemplo de pesos que satisfagan $\frac{1}{N}\sum_i w_i=1$ pero para los
              cuales la varianza \emph{no} converge a cero.
    \end{enumerate}
    Nota: si $\lim_{N\to\infty}\E(\hat\theta_N)=\theta$ y $\lim_{N\to\infty}\Var(\hat\theta_N)=0$,
    entonces $\hat\theta_N\pto\theta$.

%% ---------------------------------------------------------------- 2
\item \emph{Insesgadez y consistencia.} Indique si cada afirmación es verdadera o falsa y
    justifique (con una demostración o con un contraejemplo).
    \begin{enumerate}
        \item Un estimador insesgado es siempre consistente.
        \item Un estimador consistente es siempre insesgado.
        \item Si $\hat\theta_N\pto\theta$, entonces $\lim_{N\to\infty}\E(\hat\theta_N)=\theta$.
    \end{enumerate}
    Pistas: para (a) considere $\tilde y = y_1$; para (b) considere la varianza muestral
    $S^2=\frac{1}{N}\sum_i (y_i-\bar y)^2$; para (c) considere la secuencia $z_N$ que toma el
    valor $0$ con probabilidad $1-1/N$ y el valor $N$ con probabilidad $1/N$.

%% ---------------------------------------------------------------- 3
\item \textbf{(*)} \emph{LGN y teorema de la función continua.} Dada una muestra i.i.d.\
    $\{(y_i,x_i)\}$, encuentre el límite en probabilidad de los siguientes estadísticos e
    indique qué supuestos necesita.
    \begin{enumerate}
        \item $\displaystyle \frac{1}{N}\sum_{i=1}^N x_i\,y_i$.
        \item Considere el modelo de regresión \emph{sin constante}, con $x$ escalar:
              $\LP(y|x)=\beta\,x$, es decir $\beta=\arg\min_b \E[(y-x\,b)^2]$.
              \begin{enumerate}[label=(\roman*)]
                  \item Obtenga la condición de primer orden y demuestre que
                        $\beta=\E(xy)/\E(x^2)$. \\
                        Nota: en clase vimos la fórmula general
                        $\beta=\E(xx')^{-1}\E(xy)$ para $x=(1,x_1,\dots,x_K)'$ \emph{con}
                        constante; aquí solo hay que particularizarla al caso escalar.
                  \item Demuestre que $\displaystyle
                        \hat\beta=\frac{\frac{1}{N}\sum_{i=1}^N x_i y_i}{\frac{1}{N}\sum_{i=1}^N x_i^2}$
                        es un estimador consistente de ese $\beta$.
              \end{enumerate}
        \item Recuerde que el coeficiente de correlación poblacional entre $x$ e $y$ es
              \begin{align*}
              \Corr(x,y)=\frac{\Cov(x,y)}{\sqrt{\Var(x)\,\Var(y)}},
              \qquad \Cov(x,y)=\E(xy)-\E(x)\E(y),
              \end{align*}
              y su análogo muestral es
              \begin{align*}
              \hat\rho = \frac{S_{xy}}{\sqrt{S_x^2\,S_y^2}},
              \qquad
              S_{xy}=\frac{1}{N}\sum_i x_i y_i-\bar x\,\bar y,
              \qquad
              S_x^2=\frac{1}{N}\sum_i x_i^2-\bar x^2,
              \end{align*}
              y análogamente $S_y^2$. Demuestre que $\hat\rho\pto\Corr(x,y)$. \\
              Pista: no necesita volver a los cinco promedios muestrales. En clase ya se
              demostró que la varianza muestral es un estimador consistente de la varianza
              poblacional, $S_x^2\pto\Var(x)$ y $S_y^2\pto\Var(y)$; demuestre primero que
              $S_{xy}\pto\Cov(x,y)$ con el mismo argumento y después aplique el teorema de la
              función continua una sola vez.
    \end{enumerate}

%% ---------------------------------------------------------------- 4
\item \textbf{(*)} \emph{Convergencia en distribución, función continua y Slutsky.} Suponga que
    $z_N \dto z \sim \N(0,1)$, que $c_N \pto 3$ y que $a_N \pto 0$. Encuentre la distribución
    asintótica de
    \begin{enumerate}
        \item $x_N = 2\,z_N+5$;
        \item $x_N = z_N^2$;
        \item $x_N = z_N/c_N$;
        \item $x_N = z_N + a_N$.
    \end{enumerate}
    En cada caso indique si usa el teorema de la función continua (para $\dto$) o el teorema de
    Slutsky.

%% ---------------------------------------------------------------- 5
\item \textbf{(*)} \emph{Teorema central del límite.} Sea $\hat\mu_3=\frac{1}{N}\sum_{i=1}^N y_i^3$
    un estimador de $\mu_3=\E(y^3)$.
    \begin{enumerate}
        \item Demuestre que $\sqrt{N}(\hat\mu_3-\mu_3)\dto \N(0,v^2)$ para algún $v^2$.
        \item Escriba $v^2$ como función de los momentos de $y$.
        \item Proponga un estimador consistente $\hat v^2$ de $v^2$ y demuestre que lo es.
        \item Encuentre la distribución asintótica de $\sqrt{N}(\hat\mu_3-\mu_3)/\hat v$.
              ¿Por qué es útil este resultado?
    \end{enumerate}

%% ---------------------------------------------------------------- 6
\item \textbf{(*)} \emph{Vectores aleatorios y formas cuadráticas.} Suponga que
    $\mathbf{z}_N \dto \mathbf{z}\sim \N(\boldsymbol{\mu},\boldsymbol{\Sigma})$, donde
    $\mathbf{z}_N=(z_{1N},z_{2N})'$ es un vector de $2\times 1$,
    $\boldsymbol{\mu}=(5,-1)'$ y
    \begin{align*}
    \boldsymbol{\Sigma}=\begin{pmatrix} 4 & 0 \\ 0 & 9 \end{pmatrix}.
    \end{align*}
    \begin{enumerate}
        \item Encuentre la distribución asintótica de $\mathbf{A}\,\mathbf{z}_N$, donde
              \begin{align*}
              \mathbf{A}=\begin{pmatrix} 1 & 1 \\ 1 & -1 \end{pmatrix}.
              \end{align*}
        \item Encuentre la distribución asintótica de
              $\boldsymbol{\Sigma}^{-1/2}(\mathbf{z}_N-\boldsymbol{\mu})$, donde
              \begin{align*}
              \boldsymbol{\Sigma}^{-1/2}=\begin{pmatrix} 1/2 & 0 \\ 0 & 1/3 \end{pmatrix}.
              \end{align*}
        \item Encuentre la distribución asintótica de
              $(\mathbf{z}_N-\boldsymbol{\mu})'\boldsymbol{\Sigma}^{-1}(\mathbf{z}_N-\boldsymbol{\mu})$.
        \item Suponga que reemplazamos $\boldsymbol{\Sigma}$ por un estimador consistente
              $\hat{\boldsymbol{\Sigma}}\pto\boldsymbol{\Sigma}$. ¿Cambian los resultados
              anteriores? ¿Por qué importa esto en la práctica?
    \end{enumerate}

%% ---------------------------------------------------------------- 7
\item \textbf{(*)} \emph{Un estimador no estándar.} Dado el modelo
    \begin{align*}
    y_i = \beta\,x_i + u_i,
    \end{align*}
    donde $\E(u_i|x_i)=0$, $\beta$ es un escalar (no es un vector) y $x_i$ es una variable
    aleatoria (no es un vector aleatorio). Considere el estimador
    \begin{align*}
    \hat\beta = \frac{\frac{1}{N}\sum_{i=1}^N x_i^3 y_i}{\frac{1}{N}\sum_{i=1}^N x_i^4}.
    \end{align*}
    \begin{enumerate}
        \item Muestre que $\hat\beta$ es un estimador consistente de $\beta$.
        \item Encuentre la distribución asintótica de $\sqrt{N}(\hat\beta-\beta)$.
        \item Simplifique la varianza asintótica bajo el supuesto de que la varianza
              condicional de $u$ dado $x$ es constante, $\E(u^2|x)=\sigma^2$.
        \item Sea $\gamma=1/\beta$ con $\beta\neq 0$. Proponga un estimador consistente
              $\hat\gamma$ y obtenga la distribución asintótica de $\sqrt{N}(\hat\gamma-\gamma)$.
    \end{enumerate}

%% ---------------------------------------------------------------- 8
\item \emph{Método delta cuando la derivada se anula.} Sea $\bar y$ la media muestral de una
    muestra i.i.d.\ con $\E(y)=\mu$ y $\Var(y)=\sigma^2$, y considere $\bar y^2$ como estimador
    de $\mu^2$.
    \begin{enumerate}
        \item Aplique el método delta para mostrar que
              \begin{align*}
              \sqrt{N}\left(\bar y^2-\mu^2\right)\dto \N\left(0,\ 4\mu^2\sigma^2\right).
              \end{align*}
        \item Suponga ahora que $\mu=0$. ¿Qué varianza asintótica entrega la fórmula de la parte
              (a)? ¿Por qué el método delta deja de ser informativo en este caso?
        \item Manteniendo $\mu=0$, encuentre la distribución asintótica de
              $N\,\bar y^2/\sigma^2$. \\
              Pista: escriba $N\bar y^2=\left(\sqrt{N}\,\bar y\right)^2$, aplique el TCL a
              $\sqrt{N}\,\bar y$ y después el teorema de la función continua.
        \item Compare (a) y (c). ¿A qué velocidad converge $\bar y^2$ a $\mu^2$ en cada caso, y
              a qué familia de distribuciones pertenece el límite?
    \end{enumerate}

%% ---------------------------------------------------------------- 9
\item \emph{Proyección lineal ponderada.} Estamos interesados en estimar un parámetro $\theta$
    definido como
    \begin{align*}
    \theta = \arg\min_t\ \E\left[(y-x\,t)^2 f(x)\right],
    \end{align*}
    donde $y$ es una variable aleatoria con $\E(y)=0$, $x$ es una variable aleatoria (no es un
    vector) con $\E(x)=0$, y $f(x)$ es una función de $x$ que toma valores entre cero y uno
    según el valor de $x$. Responda enunciando los teoremas que utilice.
    \begin{enumerate}
        \item Encuentre una expresión explícita para $\theta$.
        \item Demuestre que si $\E(y|x)=\beta\,x$, entonces $\theta$ coincide con el coeficiente
              del modelo de regresión $\beta=\E(xy)/\E(x^2)$.
        \item Suponga que tiene una muestra aleatoria de tamaño $N$ de $(y,x)$. En base a su
              respuesta en (a), proponga un estimador $\tilde\theta$ y demuestre que es consistente.
        \item Encuentre la distribución asintótica de $\sqrt{N}(\tilde\theta-\theta)$. Puede
              utilizar el siguiente resultado: $\E[x\,u\,f(x)]=0$, donde $u=y-\theta\,x$.
    \end{enumerate}

\end{enumerate}

\vspace{1em}
\noindent\rule{\textwidth}{0.4pt}

\section*{Práctica computacional (opcional)}

\noindent En el repositorio del curso está el ejemplo
\href{https://github.com/ramirode/econometria_mae/tree/main/examples/3.teoria_asintotica_simulacion-bernoulli}%
{\texttt{3.teoria\_asintotica\_simulacion-bernoulli}}, que simula 50.000 muestras de una
Bernoulli($p=0.78$) para $N=2,5,25,100$ y grafica la distribución muestral de $\bar y_N$ y de
$\sqrt{N}(\hat p - p)/\sqrt{p(1-p)}$.

\begin{enumerate}[itemsep=0.6em]
    \item Corra el script y verifique visualmente la consistencia de $\bar y_N$ y el teorema
          central del límite.
    \item Modifique $p$ a un valor cercano a cero (por ejemplo $p=0.02$). ¿Qué tan grande tiene
          que ser $N$ para que la aproximación normal sea razonable? Relacione su respuesta con
          el hecho de que el TCL es un resultado asintótico.
\end{enumerate}

\end{document}
